% !TeX root = ../drafts/c-flock-protocol.tex
\section{The Flock protocol}\label{annex:flock}

Flock~\cite{Flock26} proves many executions of one Boolean circuit at once, in our case the BLAKE2s compression. Since we do not follow their protocol exactly, we describe our version here.

\subsection{Statement and commitment}\label{flock:statement}

One BLAKE2s compression has a Boolean R1CS witness block of $2^{\kblock}$ bits, where $\kblock=14$. Its fixed matrices are $A_0,B_0,C_0\in\Ftwo^{2^{\kblock}\times 2^{\kblock}}$, with $C_0=I$. The witness includes the circuit inputs, outputs, AND-gate values, addition carries and one constant-one position; linear wires are substituted away. Within-block coordinates are the low-order coordinates and batch coordinates are the high-order coordinates throughout this annex, and a Boolean tuple is also read as an integer, lowest coordinate first.

For $2^{\kbatch}$ independent compressions, the witness is:
\[
z:\cube{\kbool}\longrightarrow\Ftwo,\qquad \kbool=\kblock+\kbatch,
\]
and the batched matrices are
\[
A=I_{2^{\kbatch}}\otimes A_0,\qquad B=I_{2^{\kbatch}}\otimes B_0,\qquad C=I_{2^{\kbool}}
\]
Writing $a=Az$, $b=Bz$ and $c=Cz=z$, the R1CS statement is:
\begin{equation}\label{flock:eq:r1cs}
a(u)b(u)+c(u)=0\qquad\text{for every }u\in\cube{\kbool}
\end{equation}

We also enforce the position $512$ of every block to be $1$ (preventing the all-zero witness):
\begin{equation}\label{flock:eq:const-one}
z(0,0,0,0,0,0,0,0,0,1,0,0,0,0,\,t)=1\qquad\text{for every }t\in\cube{\kbatch}
\end{equation}

This second condition is enforced by lincheck (\S\ref{flock:lincheck}), and position $512=\jconst$ is where every circuit below puts its constant wire. Since it is asked of every block, the prover pads the batch out to $2^{\kbatch}$ with copies of a genuine witness block rather than with zeros.

Set $\kskip=6$ and $\nflock=\kbool-\kskip$. The prover commits to $z$ by packing its low $\kskip$ coordinates, $64$ bits, to a $\K$-element, following Ring-Switching (Annex~\ref{annex:rs}). The resulting multilinear $\qflock:\cube{\nflock}\to\K$ is committed as one region of the stack (\S\ref{sec:stacking}).


\subsection{Zerocheck}\label{flock:zerocheck}

Flock zerochecks \eqref{flock:eq:r1cs} (Definition~\ref{def:zerocheck}), but replaces the first $\kskip$ rounds by one round through a multiplication code~\cite{DBHR26}, the \emph{AG skip}. It is the univariate skip of~\cite{Gru24} with the Reed-Solomon code exchanged for an algebraic-geometry code, chosen so that the product of two binary messages is again binary: the prover's round-one message is then a sum of selected equality weights, with no field product per witness bit. Split $u=(i,v)$ with $i\in\cube{\kskip}$ and $v\in\cube{\nflock}$.

\paragraph{The code.} Let $\skipcurve$ be the curve over $\Ftwo$ of genus $95$ whose affine model is
\begin{equation}\label{flock:eq:curve}
y^4+(x+1)y^2+(x^3+x)y+x^7+x^3=0,\qquad z_i^2+z_i=g_i(x,y)\quad(i=1,2,3),
\end{equation}
with $g_1,g_2,g_3$ the rational functions of~\cite[Eq.~(A.1)]{DBHR26}, whose denominators divide $d(x)=(x+1)(x^{10}+x^4+1)$. It has exactly $64$ points over $\Ftwo$, the \emph{skip domain} $\skipdom=\{P_0,\ldots,P_{63}\}$, and~\cite{DBHR26} gives a divisor $D$ of degree $158$ on it such that the functions $L(D)$ with poles bounded by $D$ form a space of dimension $64$ on which evaluation at $\skipdom$ is a bijection onto $\Ftwo^{64}$. Write $\lag_i\in L(D)$ for the function with $\lag_i(P_{i'})=[i=i']$. Products of two such functions lie in $L(2D)$, of dimension $222$.

\begin{fact}[Pairs]\label{fact:skip-pairs}
There is a set $\skippairs$ of $158$ pairs $i<k$ such that the $64$ squares $\lag_i^2$ and the $158$ products $\lag_i\lag_k$, $(i,k)\in\skippairs$, form a basis of $L(2D)$. A function of $L(2D)$ vanishes on $\skipdom$ if and only if its coordinates over the squares are zero.
\end{fact}
The second sentence follows from the first: a product $\lag_i\lag_k$ with $i\neq k$ vanishes at every point of $\skipdom$, and $\lag_i^2$ at every point but $P_i$. The first is a rank computation on the extension matrix of~\cite[\S A.3]{DBHR26}, and $\skippairs$ is the greedy choice in lexicographic order.

The codewords are evaluations over the challenge field. The \emph{evaluation set} $\skipset$ is the set of solutions $(x,y,z_1,z_2,z_3)\in\E^5$ of \eqref{flock:eq:curve}, its second equation multiplied through by $d(x)$, at which $d(x)\,\eta(x)\,\eta_0(x)\neq 0$, where $\eta$ of degree $49$ and $\eta_0$ of degree $31$ are the common denominators of the bases of $L(2D)$ and $L(D)$. Every $\lag_i$ is $\eta_0(x)^{-1}$ times a polynomial in $x,y,z_1,z_2,z_3$ with coefficients in $\Ftwo$, so the verifier evaluates all $64$ at a point of $\skipset$ from the point's coordinates.

\begin{fact}[Distance~{\cite[Eq.~(B.4)]{DBHR26}}]\label{fact:skip-distance}
The zeros in $\skipset$ of a nonzero $F\in L(2D)\otimes\E$ lie over at most $2248$ values of $x$.
\end{fact}
Each such zero lies over a root of the norm of $\eta F$ down to $\E(x)$, a nonzero rational function whose numerator has degree at most $2248$. This is the bound~\cite{DBHR26} certifies from the equations alone; the geometric one, $2\deg D=316$ zeros, is not needed here.

Write $\lag_i$ also for these evaluations.

\begin{definition}[Quirky extension]\label{def:quirky}
For $f:\cube{\kskip}\times\cube{n}\to\Ftwo$, its \emph{quirky extension} extends the first $\kskip$ coordinates through the code and the other $n$ multilinearly:
\[
  \qext f:\skipset\times\E^{n}\to\E,\qquad
  \qext f(Y,v):=\sum_{i=0}^{63}\lag_i(Y)\,\mle f(i,v),
\]
so $\qext f(\cdot,v)$ lies in $L(D)\otimes\E$, $\qext f$ has degree one in every other variable, and $\qext f(P_i,v)=\mle f(i,v)$ for every $i\in\cube{\kskip}$. Below $f$ is $a$, $b$ or $c$, with $n=\nflock$, or $A_0$ or $B_0$ read as a row followed by a column, with $n=2\kblock-\kskip$. Its argument $(Y,v)$ is a \emph{quirky point}.
\end{definition}

The zerocheck point $r \in \E^{\nflock}$ is defined as follows:
\begin{itemize}
  \item Its first seven coordinates $r_0, \dots, r_6$ are fixed to $\gen^{2^j}/(1+\gen^{2^j})$ for $0\leq j<6$ and to $y/(1+y)$, with $\gen$ the generator of $\K$ and $y$ that of $\E$ over $\K$. Their $2^7$ equality weights are the elements $\gen^{k}y^{t}$, $0\leq k<64$, $t\in\{0,1\}$, divided by one constant: they are part of a basis of $\E$ over $\Ftwo$, so Lemma~\ref{lem:fixed-zerocheck} applies. The first six are what lets the prover read $64$ consecutive blocks of one coordinate as one element of $\K$.
  \item The remaining coordinates $r_7,\dots,r_{\nflock-1}$ are sampled at random by the verifier.
\end{itemize}

Define
\begin{equation}\label{flock:eq:pskip}
P(Y)=\sum_{v\in\cube{\nflock}}\eq(r,v)\bigl(\qext a(Y,v)\qext b(Y,v)+\qext c(Y,v)\bigr),
\end{equation}
which lies in $L(2D)\otimes\E$. If \eqref{flock:eq:r1cs} holds, then $P(P_i)=0$ for every $i$, so by Fact~\ref{fact:skip-pairs}
\begin{equation}\label{flock:eq:pskip-pairs}
P=\sum_{(i,k)\in\skippairs}p_{ik}\,\lag_i\lag_k
\end{equation}
and the prover sends the $158$ coefficients $p_{ik}\in\E$: the $64$ over the squares are the zeros the verifier assumes.

The verifier then samples $x_0\in\E$. About one $x$ in twelve has a point of $\skipset$ over it, so the prover answers with a counter $\delta\in\K$, an integer below $2^{16}$ read as bits, and with $y,z_1,z_2,z_3\in\E$. The verifier checks the range of $\delta$ and that
\[
\zskip=(x_0+\omega\delta,\,y,\,z_1,\,z_2,\,z_3)\in\skipset,
\]
where $\omega\in\E$ is a fixed public constant outside $\K$. A step inside $\K$ would not do: the first condition for a point to lie over $x$ is $\operatorname{Tr}(x^3+x)=0$, the trace of $\K$ vanishes on the polynomials in $\gen$ of degree below $60$, and a counter stepping by $1$ would then move that trace only linearly, leaving one $x_0$ in $2^{17}$ with no counter that passes. The verifier
then evaluates the $64$ weights $\lag_i(\zskip)$, and reconstructs the claimed value
\[
v_P \qeq \sum_{(i,k)\in\skippairs}p_{ik}\,\lag_i(\zskip)\,\lag_k(\zskip).
\]

The parties run sumcheck over the remaining $\nflock$ Boolean variables in \eqref{flock:eq:pskip}, with Gruen's eq-factor optimization~\cite{Gru24}. Writing $R_{\mathrm{zc}}\in\E$ for the running claim after the last round, the prover sends $v_a,v_b\in\E$ and the terminal identity
\begin{equation}\label{flock:eq:zc-terminal}
R_{\mathrm{zc}}=v_av_b+v_c
\end{equation}
is what \emph{defines} $v_c$: both sides solve it rather than transmit it. The zerocheck therefore ends leaving three claims at one point,
\begin{equation}\label{flock:eq:zc-claims}
\qext a(\zskip,\fc)=v_a,\qquad \qext b(\zskip,\fc)=v_b,\qquad \qext z(\zskip,\fc)=v_c,
\end{equation}
where $\fc\in\E^{\nflock}$ is the vector of sumcheck challenges and the third claim uses $c=z$.

\subsection{Lincheck}\label{flock:lincheck}

Lincheck reduces the three zerocheck claims to claims on $z$. The block structure is what keeps it cheap: the matrices repeat one BLAKE2s circuit, so lincheck runs over a single block, not over the batch. Split $\fc=(\fc_{\mathrm{in}},\fc_{\mathrm{out}})$, where $\fc_{\mathrm{in}}\in\E^{\kblock-\kskip}$ selects the remaining within-block coordinates and $\fc_{\mathrm{out}}\in\E^{\kbatch}$ selects a batch instance.

\begin{lemma}[Block diagonality folds the batch into $z$]\label{lem:flock-block-diagonal}
For $a=Az$ with $A=I_{2^{\kbatch}}\otimes A_0$, at every point $(\zskip,\fc_{\mathrm{in}},\fc_{\mathrm{out}})\in\skipset\times\E^{\kblock-\kskip}\times\E^{\kbatch}$,
\[
\qext a(\zskip,\fc_{\mathrm{in}},\fc_{\mathrm{out}})=\sum_{j\in\cube{\kblock}}\qext A_0(\zskip,\fc_{\mathrm{in}},j)\,\mle z(j,\fc_{\mathrm{out}}),
\]
and likewise for $b=Bz$ with $B_0$ in place of $A_0$.
\end{lemma}

\begin{proof}
Block diagonality gives
\[
a(k,t)=\sum_{j\in\cube{\kblock}}A_0(k,j)\,z(j,t)\qquad\text{for }k\in\cube{\kblock},\;t\in\cube{\kbatch},
\]
so, splitting the row as $k=(\ksk,\kin)$ with $\ksk\in\cube{\kskip}$ and $\kin\in\cube{\kblock-\kskip}$ and expanding every extension by Fact~\ref{fact:eq},
\begin{align*}
\qext a(\zskip,\fc_{\mathrm{in}},\fc_{\mathrm{out}})&=\sum_{\ksk\in\cube{\kskip}}\lag_{\ksk}(\zskip)\,\mle a(\ksk,\fc_{\mathrm{in}},\fc_{\mathrm{out}})\\
&=\sum_{\ksk}\lag_{\ksk}(\zskip)\sum_{\kin}\eq(\fc_{\mathrm{in}},\kin)\sum_{t\in\cube{\kbatch}}\eq(\fc_{\mathrm{out}},t)\,a\bigl((\ksk,\kin),t\bigr)\\
&=\sum_{\ksk,\kin}\lag_{\ksk}(\zskip)\eq(\fc_{\mathrm{in}},\kin)\sum_{j\in\cube{\kblock}}A_0((\ksk,\kin),j)\underbrace{\sum_{t}\eq(\fc_{\mathrm{out}},t)\,z(j,t)}_{=\;\mle z(j,\fc_{\mathrm{out}})}\\
&=\sum_{j\in\cube{\kblock}}\Bigl(\underbrace{\sum_{\ksk,\kin}\lag_{\ksk}(\zskip)\,\eq(\fc_{\mathrm{in}},\kin)\,A_0((\ksk,\kin),j)}_{=\;\qext A_0(\zskip,\fc_{\mathrm{in}},j)}\Bigr)\mle z(j,\fc_{\mathrm{out}}).
\end{align*}
\end{proof}

The $c$ claim has the same shape, with the column index split as $j=(\jsk,\jin) \in\cube{\kskip} \times \cube{\kblock-\kskip}$, $\jsk$ the coordinates the AG skip replaced and $\jin$ the within-block coordinates that remain:
\begin{equation}\label{flock:eq:c-as-column}
\qext z(\zskip,\fc_{\mathrm{in}},\fc_{\mathrm{out}})=\sum_{j\in\cube{\kblock}}\lag_{\jsk}(\zskip)\,\eq(\fc_{\mathrm{in}},\jin)\,\mle z(j,\fc_{\mathrm{out}}),
\end{equation}

\vspace{3mm}

The verifier samples $\lcalpha\in\E$ and batches the zerocheck's three claims \eqref{flock:eq:zc-claims} with the constant-position check \eqref{flock:eq:const-one}:
\begin{equation}\label{flock:eq:lincheck}
\begin{split}
v_a+\lcalpha v_b+\lcalpha^2v_c+\lcalpha^3\qeq\sum_{j\in\cube{\kblock}}\Bigl(&\qext A_0(\zskip,\fc_{\mathrm{in}},j)+\lcalpha\qext B_0(\zskip,\fc_{\mathrm{in}},j)\\
&+\lcalpha^2\lag_{\jsk}(\zskip)\eq(\fc_{\mathrm{in}},\jin)+\lcalpha^3[j=512]\Bigr)\mle z(j,\fc_{\mathrm{out}}).
\end{split}
\end{equation}

\vspace{2mm}

The parties start by running $\kblock - \kskip = 14 - 6 = 8$ sumcheck rounds on $\eqref{flock:eq:lincheck}$, binding the $\jin$ coordinates (leaving the skip index $\jsk$ unfolded), with challenges $\fc'_{\mathrm{in}}$ in coordinate order (the reverse of their round order). Let $R_{\mathrm{lc}}\in\E$ be the resulting sumcheck claim.

The prover sends $(s_0,\ldots,s_{63})\in\E^{64}$, claiming $s_i=\mle z(i,\fc'_{\mathrm{in}},\fc_{\mathrm{out}})$. These 64 claims will later be proven by Ring-Switching.

The verifier finally needs to check:

\begin{equation}\label{flock:eq:lincheck-terminal}
\begin{split}
R_{\mathrm{lc}}\qeq\sum_{i=0}^{63}\Bigl(&\qext A_0(\zskip,\fc_{\mathrm{in}},i,\fc'_{\mathrm{in}})+\lcalpha\,\qext B_0(\zskip,\fc_{\mathrm{in}},i,\fc'_{\mathrm{in}})\\
&+\lcalpha^2\lag_i(\zskip)\,\eq(\fc_{\mathrm{in}},\fc'_{\mathrm{in}})+\lcalpha^3[i=0]\,\eq(\fc'_{\mathrm{in}},8)\Bigr)s_i.
\end{split}
\end{equation}

\subsubsection{Concluding Lincheck for a native verifier}\label{flock:native}


With the row split $k=(\ksk,\kin) \in\cube{\kskip} \times \cube{\kblock-\kskip}$ of the proof above,
\begin{equation}\label{flock:eq:erow}
\qext A_0(\zskip,\fc_{\mathrm{in}},j)=\sum_{k\in\cube{\kblock}}\erow(k)\,A_0(k,j),\qquad
\erow(\ksk,\kin)=\lag_{\ksk}(\zskip)\,\eq(\fc_{\mathrm{in}},\kin),
\end{equation}
and likewise for $B_0$.

Writing $\wcol(\jsk,\jin)=s_{\jsk}\,\eq(\fc'_{\mathrm{in}},\jin)$ with $(\jsk,\jin) \in\cube{\kskip} \times \cube{\kblock-\kskip}$, each matrix term in $\eqref{flock:eq:lincheck-terminal}$ is simplified as follows:
\begin{equation}\label{flock:eq:matrix-form}
\sum_{i=0}^{63}\qext A_0(\zskip,\fc_{\mathrm{in}},i,\fc'_{\mathrm{in}})\,s_i=\sum_{k,j\in\cube{\kblock}}A_0(k,j)\,\erow(k)\,\wcol(j).
\end{equation}

One forward walk of the BLAKE2s circuit returns \eqref{flock:eq:matrix-form} for both matrices (\S\ref{flock:matrices}), never touching the roughly $89$ million nonzero entries they hold. Its committed wires are $896$ free inputs, the constant, $256$ committed $\oplus$ wires carrying the output chaining value, and $14{,}720$ $\wedge$ wires, one per product bit of the $320$ modular additions of the ten rounds; at one multiplication per committed wire on the $A$ side, one per $\wedge$ wire on the $B$ side, and one for the constant factor the remaining $B$ rows share, that is $30{,}594$ field multiplications in $O(\text{circuit})$ work. The $c$ term of \eqref{flock:eq:lincheck-terminal} is $\lcalpha^2\inner{\erow}{\wcol}$ in the same notation, and needs no walk: both sides are tensors, so it contracts to $\lcalpha^2\eq(\fc_{\mathrm{in}},\fc'_{\mathrm{in}})\sum_i\lag_i(\zskip)s_i$.

\subsubsection{Concluding Lincheck for a recursive verifier}\label{flock:defer}

To avoid the costly circuit-walk in recursion, the recursion program receives as hint:
\begin{equation}\label{flock:eq:matpart}
M_{\mathrm{lc}}=\sum_{k,j\in\cube{\kblock}}\bigl(A_0(k,j)+\lcalpha B_0(k,j)\bigr)\erow(k)\,\wcol(j).
\end{equation}
This makes it possible to cheaply verify \eqref{flock:eq:lincheck-terminal}. Ensuring validity of the hint is deferred outside of the recursion program, via sumcheck reduction, as described in \S\ref{sec:deferred-claims}.

\subsubsection{Ring switching}\label{flock:ringswitch}

The role of Ring-Switching is now to prove:
\[
s_i \qeq \mle z(i,\fc'_{\mathrm{in}},\fc_{\mathrm{out}})\qquad\text{for }i\in\cube{\kskip}.
\]


\subsection{Protocol summary}\label{flock:summary}

\begin{protocol}[Flock summary]\label{flock:protocol}
After the witness $z$ has been committed, and its size announced:
\begin{enumerate}
\item The verifier forms the zerocheck point $r$: seven fixed coordinates, the rest sampled. The $\kskip$ variables the AG skip replaced consume no equality challenge.
\item The prover sends the $158$ coefficients of $P$ over the pairs \eqref{flock:eq:pskip-pairs}. The verifier samples $x_0$, the prover answers with the counter $\delta$ and the rest of the point $\zskip$, and the verifier checks $\zskip\in\skipset$ and reconstructs $v_P$.
\item The parties run $\nflock$ quadratic zerocheck rounds on \eqref{flock:eq:pskip}, which the prover ends by sending $v_a,v_b$; the terminal identity \eqref{flock:eq:zc-terminal} then fixes $v_c$.
\item The verifier samples $\lcalpha$ and the parties run the $8$ quadratic lincheck rounds on \eqref{flock:eq:lincheck}. The prover sends $(s_0,\ldots,s_{63})$, and the verifier checks \eqref{flock:eq:lincheck-terminal}: one circuit walk natively, the hint \eqref{flock:eq:matpart} in recursion.
\end{enumerate}
What Flock leaves is one family of $64$ claims for Ring-Switching (\S\ref{flock:ringswitch}): $(s_i \qeq \mle z(i, \fc'_{\mathrm{in}},\fc_{\mathrm{out}}))_{0 \leq i \leq 63}$.
\end{protocol}

\subsection{Evaluating the matrices}\label{flock:matrices}

$A_0$ and $B_0$ encode a Boolean circuit, but are never materialized.

\subsubsection{The circuit}\label{flock:circuit}

\begin{definition}[Circuit]\label{def:circuit}
A \emph{circuit} is an ordered sequence of wires $w_1,\ldots,w_N$. Each wire is either a \emph{source}, a $\oplus$ wire with a possibly empty set $S_w$ of earlier wires, or a $\wedge$ wire with an ordered pair $(u_w,u'_w)$ of earlier wires. Its value lies in $\Ftwo$: a $\oplus$ wire is the sum of the values in $S_w$, and a $\wedge$ wire is the product of the values of $u_w$ and $u'_w$. The source $w_1$ is the constant $1$ and the other sources are free inputs.

A set of \emph{committed} wires containing all sources, all $\wedge$ wires, and a subset of the $\oplus$ wires (typically the result of the computation) is equipped with an injection $w\mapsto\posof{w}$ into $\cube{\kblock}$. The constant has position $\posof{w_1}=\jconst$. A position $k\in\cube{\kblock}$ is \emph{empty} if $k\neq\posof{w}$ for every committed wire $w$. A witness block $z_t=z(\cdot,t)$ stores each committed wire's value at its position and stores $0$ at empty positions.
\end{definition}

For $p\in\cube{\kblock}$, write $[\,\cdot=p\,]:\cube{\kblock}\to\Ftwo$ for the function that equals $1$ at $p$ and $0$ at every other position. The \emph{expansion} of a wire over committed positions is
\begin{equation}\label{flock:eq:expand}
\lform{w}=[\,\cdot=\posof{w}\,]\ \text{ if $w$ is committed},\qquad \lform{w}=\sum_{u\in S_w}\lform{u}\ \text{ otherwise},
\end{equation}
where the second case applies only to uncommitted $\oplus$ wires. By induction, $\inner{\lform{w}}{z_t}$ is the value of $w$.

\subsubsection{The R1CS matrices}\label{flock:rows}

For every committed wire $w$, assign two sets of wires to the corresponding row:
\[
(S^A_w,S^B_w)=
\begin{cases}
(\{u_w\},\{u'_w\}) & \text{if $w$ is a $\wedge$ wire},\\
(S_w,\{w_1\}) & \text{if $w$ is a $\oplus$ wire},\\
(\{w\},\{w_1\}) & \text{if $w$ is a source}.
\end{cases}
\]
and define:
\begin{equation}\label{flock:eq:row-lists}
A_0(\posof{w},\cdot)=\sum_{u\in S^A_w}\lform{u},\qquad B_0(\posof{w},\cdot)=\sum_{u\in S^B_w}\lform{u}.
\end{equation}
For every empty position $k$, set $A_0(k,\cdot)=B_0(k,\cdot)=0$. This specifies every row of both matrices.

\vspace{3mm}

The R1CS constraints given by $A_0$, $B_0$ and $C_0=I$ faithfully encode the circuit.
\subsubsection{Forward algorithm}\label{flock:forward}

Given a column vector $W:\cube{\kblock}\to\E$, the forward algorithm computes the two row-indexed vectors $A_0W$ and $B_0W$ without constructing either matrix. The native lincheck verifier uses $W=\wcol$ and retains only the scalar $\inner{\erow}{A_0W}+\lcalpha\inner{\erow}{B_0W}$. The prover of the deferred matrix sumcheck in \S\ref{sec:deferred-claims} instead stores both vectors and folds them through its first $\kblock$ rounds.

To compute them, visit the wires from $w_1$ to $w_N$ and evaluate every expansion against $W$:
\[
\inner{\lform{w}}{W}=W(\posof{w})\ \text{ if $w$ is committed},\qquad \inner{\lform{w}}{W}=\sum_{u\in S_w}\inner{\lform{u}}{W}\ \text{ otherwise};
\]
then \eqref{flock:eq:row-lists} gives
\begin{equation}\label{flock:eq:row-pair}
\inner{A_0(\posof{w},\cdot)}{W}=\sum_{u\in S^A_w}\inner{\lform{u}}{W},\qquad \inner{B_0(\posof{w},\cdot)}{W}=\sum_{u\in S^B_w}\inner{\lform{u}}{W}.
\end{equation}
Together with zero at every empty row, these values are the desired vectors.

\subsubsection{Transpose algorithm}\label{flock:backward}

For $X\in\{A,B\}$, the transpose algorithm computes the complete column-indexed vector $\colmar^{X}=X_0^{\!\top}\erow$ without constructing $X_0$. The prover needs all $2^{\kblock}$ entries of both vectors:
\begin{equation}\label{flock:eq:marginal}
\colmar^{X}(j)=\qext X_0(\zskip,\fc_{\mathrm{in}},j)=\sum_{k\in\cube{\kblock}}X_0(k,j)\erow(k)=\bigl(X_0^{\!\top}\erow\bigr)(j),\qquad X\in\{A,B\},
\end{equation}
where the second equality uses \eqref{flock:eq:erow}. Viewing each row as a column-indexed vector, empty rows contribute zero, so \eqref{flock:eq:row-lists} gives
\begin{align*}
\colmar^{X}
&=\sum_{w\ \mathrm{committed}}\erow(\posof{w})\sum_{u\in S^X_w}\lform{u}\\
&=\sum_u\left(\sum_{\substack{w\ \mathrm{committed}\\u\in S^X_w}}\erow(\posof{w})\right)\lform{u}.
\end{align*}
The last equality exchanges the finite sums and collects every occurrence of the same expansion $\lform{u}$. Thus, for every circuit wire $u$, define
\begin{equation*}
\chg{u}^{X}=\sum_{\substack{w\ \mathrm{committed}\\u\in S^X_w}}\erow(\posof{w}),
\qquad \text{so that}\qquad
\colmar^{X}=\sum_u\chg{u}^{X}\lform{u}.
\end{equation*}
Run the following procedure independently for $X=A$ and $X=B$. Initialize the output table $\colmar^{X}$ to zero and visit $w_N,\ldots,w_1$. At wire $w$, eliminate the term $\chg{w}^{X}\lform{w}$ using the appropriate case of \eqref{flock:eq:expand}:
\begin{itemize}
\item if $w$ is committed, then $\chg{w}^{X}\lform{w}=\chg{w}^{X}[\,\cdot=\posof{w}\,]$, so add $\chg{w}^{X}$ to $\colmar^{X}(\posof{w})$;
\item otherwise, $w$ is an uncommitted $\oplus$ wire and
\[
\chg{w}^{X}\lform{w}=\chg{w}^{X}\sum_{u\in S_w}\lform{u}=\sum_{u\in S_w}\chg{w}^{X}\lform{u},
\]
so add $\chg{w}^{X}$ to the existing coefficient $\chg{u}^{X}$ of every predecessor $u\in S_w$.
\end{itemize}

\subsection{Soundness analysis}\label{flock:soundness}


The commitment determines $\qflock$, hence $z$, before any challenge is drawn, and the PCS proves the claims opened on $\qflock$. So take the one transmitted family to be true of $z$,
\[
s_i=\mle z(i,\fc'_{\mathrm{in}},\fc_{\mathrm{out}}),\qquad i\in\cube{\kskip},
\]
up to the PCS error and the $2^{-160}$ of Annex~\ref{annex:rs}, and let us show a $z$ failing \eqref{flock:eq:r1cs} or \eqref{flock:eq:const-one} is rejected.

Lincheck comes first, and it is the only thing that checks anything: the zerocheck's own terminal identity defines $v_c$ instead of testing it, so every lie upstream reaches lincheck as a wrong triple $(v_a,v_b,v_c)$. Its terminal check \eqref{flock:eq:lincheck-terminal} is evaluated from the true $s$ and the matrices themselves, by the walk of \S\ref{flock:native} natively and, in recursion, once \eqref{flock:eq:matpart} is discharged (\S\ref{sec:deferred-claims}), so a false \eqref{flock:eq:lincheck} survives the $8$ rounds with probability at most $16/|\E|$ (Fact~\ref{fact:sumcheck}, degree two in each round). And by Lemma~\ref{lem:flock-block-diagonal} and \eqref{flock:eq:c-as-column}, \eqref{flock:eq:lincheck} says
\[
\bigl(v_a+\qext a(\zskip,\fc)\bigr)+\lcalpha\bigl(v_b+\qext b(\zskip,\fc)\bigr)+\lcalpha^2\bigl(v_c+\qext z(\zskip,\fc)\bigr)+\lcalpha^3\bigl(1+\mle z(512,\fc_{\mathrm{out}})\bigr)=0,
\]
degree three in $\lcalpha$, all four coefficients fixed before $\lcalpha$ is drawn. Acceptance therefore pins them all, except with $3/|\E|$ (Lemma~\ref{lem:sz}): $v_a$, $v_b$ and $v_c$ are the true evaluations, and $\mle z(512,\fc_{\mathrm{out}})=1$, which is \eqref{flock:eq:const-one} except with $\kbatch/|\E|$ more, $\fc_{\mathrm{out}}$ being random.

With all three pinned, the zerocheck's terminal identity \eqref{flock:eq:zc-terminal} becomes a real check again: it fixes $R_{\mathrm{zc}}$ to the true value of the summand at $\fc$.

Now let $z$ violate \eqref{flock:eq:r1cs} at $u=(i,v)$. Then $P(P_i)$ is the extension at $r$ of a nonzero Boolean multilinear (Definition~\ref{def:quirky}, Fact~\ref{fact:eq}), and the seven fixed coordinates of $r$ carry $\Ftwo$-independent weights, so it stays nonzero except with $(\nflock-7)/|\E|$ (Lemma~\ref{lem:fixed-zerocheck}): the $64$ zeros the verifier assumes on $\skipdom$ are false. Let $\hat P$ be the function \eqref{flock:eq:pskip-pairs} of the received coefficients. It vanishes on $\skipdom$ while $P$ does not, so $\hat P+P$ is a nonzero element of $L(2D)\otimes\E$, and both are fixed before $x_0$ is drawn. By Fact~\ref{fact:skip-distance}, $v_P=P(\zskip)$ then needs $x_0+\omega\delta$ to be one of at most $2248$ values. The prover chooses $\delta$ among $2^{16}$ and the point over $x_0+\omega\delta$ freely, and neither choice helps beyond that: $v_P\neq P(\zskip)$ except with $2^{16}\cdot 2248/|\E|$. The $\nflock$ zerocheck rounds then run on a false initial claim with a true terminal value, and reject except with $2\nflock/|\E|$.

Summing those terms, such a $z$ is accepted with probability at most
\[
\frac{16+3+\kbatch+(\nflock-7)+2^{16}\cdot 2248+2\nflock}{|\E|}=\frac{4\kbatch+36+2^{16}\cdot 2248}{|\E|}<2^{-164},
\]
on top of the two errors granted above.
